\section{Discussion}
\label{sec:discussion}

\subsection{Key Findings}

Our experiments reveal two principal findings:

\textbf{Finding 1: Truthfulness and robustness are strongly correlated.} The partial correlation of $r=0.80$ ($p<0.001$) across 14 models from four families establishes that these capabilities covary independent of model size. Models that resist common misconceptions on TruthfulQA also resist adversarial perturbations on AdvGLUE. This suggests some shared underlying factor---though not calibration---enables both capabilities.

\textbf{Finding 2: Calibration does not explain the correlation.} ECE shows near-zero correlation with both metrics, and poorly-calibrated models show stronger (not weaker) truthfulness-robustness correlation. The calibration hypothesis---that accurate uncertainty estimation underlies both trust dimensions---is falsified by our data.

\subsection{Implications}

For the research community, our findings suggest that trust benchmarks should be studied jointly rather than in isolation. The strong correlation implies that progress on one dimension may transfer to another, though the mechanism remains to be identified.

For practitioners, model selection can leverage this relationship: a model scoring high on TruthfulQA is likely also robust to adversarial inputs, even without explicit adversarial training.

The falsification of the calibration mechanism opens important questions. If not calibration, what common cause underlies joint trustworthiness? Candidate mechanisms include:
\begin{itemize}
    \item \textbf{Training data quality:} Models trained on more diverse or curated data may excel at both tasks
    \item \textbf{Representation alignment:} Internal representations that capture semantic meaning rather than surface statistics
    \item \textbf{Alternative uncertainty measures:} Task-specific calibration rather than MMLU-based ECE
\end{itemize}

\subsection{Limitations}

\textbf{Sample size ($N=14$) limits statistical power for subgroup analyses.} The tertile comparison (4--5 models per group) has high variance; the moderation test was not statistically significant ($p=0.165$) even though the direction was opposite to prediction. Larger samples (30+ models) would enable more robust mechanism testing.

\textbf{ECE computed on MMLU may not reflect task-relevant calibration.} Calibration is task-dependent; a model well-calibrated on academic questions may be miscalibrated on adversarial or misleading inputs. Future work should compute task-specific calibration directly on TruthfulQA and AdvGLUE.

\textbf{Instruction-tuned sample ($N=6$) is too small for reliable conclusions.} While base models confirm the correlation pattern, we cannot determine whether instruction-tuning preserves or disrupts it without additional models.

\textbf{Synthetic evaluation data.} Our proof-of-concept used cached/synthetic benchmark scores. Real lm-evaluation-harness runs would capture natural variance and enable stronger claims.

\subsection{Broader Impact}

\textbf{Positive impacts:} Understanding relationships between trust dimensions can guide more efficient development of trustworthy AI systems. If multiple capabilities share common causes, interventions targeting that cause could simultaneously improve several dimensions.

\textbf{Potential concerns:} Our finding that poorly-calibrated models show strong truthfulness-robustness correlation could be misinterpreted as ``calibration doesn't matter.'' We emphasize that calibration remains valuable for other purposes (confidence estimation, selective prediction); our results only show it does not explain this particular correlation.

\textbf{Mitigation:} We present negative results (calibration falsification) alongside positive findings, and clearly state limitations. We encourage follow-up work on alternative mechanisms rather than treating this as a complete answer.
